% !TEX root = ../discretemath.tex

\section{Неподвижные точки у случайной перестановки}

\subsection{Постановка задачи}

\begin{Definition}{Неподвижная точка и беспорядок}{}
	Пусть $\sigma\in S_n$. Индекс $i\in\{1,\dots,n\}$ — \emph{неподвижная точка}, если $\sigma(i)=i$. Перестановка без неподвижных точек (то есть $\sigma(i)\ne i$ для всех $i$) называется \emph{беспорядком} (derangement).
\end{Definition}

Перестановку удобно изображать стрелками $i\to\sigma(i)$: неподвижная точка — это петля.

\begin{figure}[H]
	\centering
	\begin{tikzpicture}[
		>=Stealth, thick, font=\small,
		v/.style={circle, draw, fill=black!80, inner sep=0pt, minimum size=6pt},
		lbl/.style={font=\small}
		]
		\foreach \i in {1,...,5}{
			\node[v,label={[lbl]below:$\i$}] (n\i) at (1.4*\i,0) {};
		}
		% sigma = (1)(2 4)(3 5): fixed point 1
		\draw[red!75!black,->] (n1) to[out=120,in=60,looseness=10] (n1);
		\node[red!75!black,lbl] at (1.4,0.95) {неподв.};
		\draw[->] (n2) to[out=55,in=125] (n4);
		\draw[->] (n4) to[out=235,in=305] (n2);
		\draw[->] (n3) to[out=55,in=125] (n5);
		\draw[->] (n5) to[out=235,in=305] (n3);
	\end{tikzpicture}
	\caption{Перестановка $\sigma=(1)(2\,4)(3\,5)$ как набор стрелок $i\to\sigma(i)$. Неподвижная точка $1$ изображается петлёй; остальные элементы входят в циклы длины~$\ge 2$ и неподвижными не являются.}
\end{figure}

\subsection{Вероятность отсутствия неподвижных точек}

Все $n!$ перестановок равновероятны, $P_r(\sigma)=1/n!$. Для каждого $i$ введём событие
\[
B_i=\{\sigma\in S_n \mid \sigma(i)=i\}.
\]
Событие «нет неподвижных точек» — это $\overline{B_1}\cap\cdots\cap\overline{B_n}$, поэтому
\[
P_r(\text{нет неподвижных точек})=1-P_r(B_1\cup\cdots\cup B_n).
\]

\subsubsection*{Применение формулы включений--исключений}

Напомним формулу включений--исключений:
\[
P_r\!\Bigl(\bigcup_{i=1}^n B_i\Bigr)
=\sum_{\varnothing\neq I\subseteq\{1,\dots,n\}}(-1)^{|I|-1}\,P_r\!\Bigl(\bigcap_{i\in I}B_i\Bigr).
\]
Отсюда
\[
P_r(\text{нет неподвижных точек})
=1+\sum_{\varnothing\neq I\subseteq\{1,\dots,n\}}(-1)^{|I|}\,P_r\!\Bigl(\bigcap_{i\in I}B_i\Bigr).
\]

Вычислим $P_r\bigl(\bigcap_{i\in I}B_i\bigr)$. Если все точки множества $I$ зафиксированы ($\sigma(i)=i$ для $i\in I$), то остальные $n-|I|$ элементов переставляются произвольно — таких перестановок $(n-|I|)!$. Значит,
\[
P_r\!\Bigl(\bigcap_{i\in I}B_i\Bigr)=\frac{(n-|I|)!}{n!}.
\]

Сгруппируем по $k=|I|$ (подмножеств мощности $k$ ровно $\binom{n}{k}$):
\[
P_r(\text{нет неподвижных точек})
=1+\sum_{k=1}^n\binom{n}{k}(-1)^k\frac{(n-k)!}{n!}.
\]
Поскольку
\[
\binom{n}{k}\frac{(n-k)!}{n!}
=\frac{n!}{k!(n-k)!}\cdot\frac{(n-k)!}{n!}
=\frac1{k!},
\]
получаем итоговую формулу
\[
\boxed{\;P_r(\text{нет неподвижных точек})
=\sum_{k=0}^n\frac{(-1)^k}{k!}
=1-\frac1{1!}+\frac1{2!}-\frac1{3!}+\cdots+\frac{(-1)^n}{n!}.\;}
\]

\begin{Remark}{}{}
	Это частичная сумма ряда для $e^{-1}$, поэтому
	\[
	P_r(\sigma\text{ не имеет неподвижных точек})\xrightarrow[n\to\infty]{}\frac1e\approx0{,}3679.
	\]
	Сходимость очень быстрая: уже при $n=6$ значение отличается от $1/e$ менее чем на $10^{-3}$.
\end{Remark}

\begin{center}
	\begin{tabular}{c c c}
		\toprule
		$n$ & $D_n=\#\{\text{беспорядков}\}$ & $p_n=D_n/n!$ \\
		\midrule
		$1$ & $0$   & $0{,}0000$ \\
		$2$ & $1$   & $0{,}5000$ \\
		$3$ & $2$   & $0{,}3333$ \\
		$4$ & $9$   & $0{,}3750$ \\
		$5$ & $44$  & $0{,}3667$ \\
		$6$ & $265$ & $0{,}3681$ \\
		\midrule
		$\infty$ & --- & $1/e\approx0{,}3679$ \\
		\bottomrule
	\end{tabular}
\end{center}

\subsection{Число беспорядков}

\begin{Definition}{Число беспорядков}{}
	$D_n$ — число перестановок из $S_n$ без неподвижных точек. Из формулы выше
	\[
	D_n=n!\sum_{k=0}^n\frac{(-1)^k}{k!}.
	\]
\end{Definition}

\begin{Remark}{}{}
	Так как $\bigl|D_n-n!/e\bigr|<\tfrac12$ при $n\ge1$, число $D_n$ есть \emph{ближайшее целое} к $n!/e$. Например, $D_4=9$ (а $4!/e\approx8{,}83$), $D_5=44$ (а $5!/e\approx44{,}15$).
\end{Remark}

\subsection{Математическое ожидание числа неподвижных точек}

Пусть $X(\sigma)$ — число неподвижных точек перестановки $\sigma$. Представим $X$ суммой индикаторов:
\[
X=\sum_{i=1}^n X_i,\qquad X_i=\mathbf{1}[\sigma(i)=i].
\]
Для каждого $i$
\[
\mathbb{E}[X_i]=P_r(\sigma(i)=i)=\frac{(n-1)!}{n!}=\frac1n.
\]
По линейности математического ожидания
\[
\mathbb{E}[X]=\sum_{i=1}^n\mathbb{E}[X_i]=n\cdot\frac1n=1.
\]

\begin{Remark}{}{}
	Замечательно, что среднее число неподвижных точек равно $1$ при \emph{любом} $n$. При этом вероятность \emph{полного} отсутствия неподвижных точек стремится к $1/e$: «в среднем одна», но довольно часто их нет вовсе. (Здесь используется линейность математического ожидания — она верна без всякой независимости индикаторов $X_i$.)
\end{Remark}
